% claim-ref: kalanchoe-desert#callus-duration
% claim-ref: kalanchoe-desert#callus-preparation
% claim-ref: kalanchoe-desert#failure-cue-scorch
% claim-ref: kalanchoe-desert#failure-cue-shrivel
% claim-ref: kalanchoe-desert#failure-cue-soft-dark-base
% claim-ref: kalanchoe-desert#insert-cutting
% claim-ref: kalanchoe-desert#plantlet-alternate-scope
% claim-ref: kalanchoe-desert#prepare-tools-and-medium
% claim-ref: kalanchoe-desert#primary-method
% claim-ref: kalanchoe-desert#root-check
% claim-ref: kalanchoe-desert#root-initiation-time
% claim-ref: kalanchoe-desert#rooting-light
% claim-ref: kalanchoe-desert#rooting-medium
% claim-ref: kalanchoe-desert#rooting-moisture
% claim-ref: kalanchoe-desert#rooting-temperature
% claim-ref: kalanchoe-desert#stem-cutting-length
% claim-ref: kalanchoe-desert#strip-lower-stem
% claim-ref: kalanchoe-desert#transfer-and-acclimate
% claim-ref: kalanchoe-desert#transplant-readiness
\CompanionHeader{PROPAGATION}{Kalanchoe 'Desert Surprise'}{Candidate Kalanchoe humilis 'Desert Surprise'}{Outdoor grow bag; cultivar identity and actual winter protection remain unverified.}
\Context{Stem cutting / spring or summer. General guidance; observe parent readiness before acting.}
\begin{minipage}[t]{.625\linewidth}\raggedright\vspace{0pt}
\Step{1}{Prepare}{Clean sharp pruners; fill a small draining pot with cactus potting soil or perlite. This is temporary rooting medium, not the established-plant mix. \textbf{[3]}}
\Step{2}{Select \& cut}{Take a healthy 3-6 in stem tip. This extension range is general succulent guidance, not a cultivar-specific requirement. \textbf{[1,3]}}
\Step{3}{Expose the lower stem}{Remove leaves from the lower half. Retain healthy upper growth and keep the cut clean. \textbf{[3]}}
\Step{4}{Dry \& callus}{Rest dry on a tray or in an empty cup for several days. Proceed by a dry, callused surface, not by date alone. \textbf{[3]}}
\Step{5}{Place upright}{Insert the callused lower stem upright in draining medium. Avoid water-glass rooting; no numeric insertion depth is established. \textbf{[3]}}
\Step{6}{Maintain conditions}{Avoid saturation and standing water. Use bright light without abrupt midday sun. The proposed 50-70 F (10-21 C; C rounded) range adapts general indoor succulent guidance. \textit{Proposed.} \textbf{[2,3,6]}}
\Step{7}{Inspect without uprooting}{Check very gently for anchoring resistance; do not repeatedly lift the cutting. Watch for sustained new growth as well as anchoring. No dependable cultivar-specific root-initiation interval was found. \textit{Proposed.} \textbf{[1,2,3]}}
\Step{8}{Transfer \& acclimate}{Move rooted material to a draining established-plant container. Start outdoor acclimation in semi-shade and increase sun gradually; protect against unsuitable cold. Use observation, not a fixed transplant day. \textbf{[1,4,6]}}
\end{minipage}\hfill\begin{minipage}[t]{.345\linewidth}\raggedright\vspace{0pt}
\Note{Gather}{Clean pruners, one small draining pot, a tray, and enough cactus medium or perlite to fill the pot.}
\Note{Milestones, not deadlines}{\begin{tabular}{@{}p{.30\linewidth}p{.60\linewidth}@{}}\textbf{Callus} & Several days; dry surface.\\\textbf{Roots} & Initiation interval unknown.\\\textbf{Transfer} & Anchored roots and sustained growth; proposed.\\\end{tabular}\par \textit{Proposed.} \textbf{[1,3]}}
\Note{Plantlet option}{Only if a confirmed candidate actually produces viable attached plantlets. RHS lists their removal; do not infer this behavior from unrelated Kalanchoe species. \textbf{[1]}}
\Note{Check 1}{Soft/dark base: check saturation, drainage and callusing; restart from sound tissue. \textit{Proposed.} \textbf{[2,3]}}
\Note{Check 2}{Shriveling: check overheating or bone-dry medium; correct without saturation. \textit{Proposed.} \textbf{[2,6]}}
\Note{Check 3}{Scorch after transfer: return to semi-shade; resume gradual acclimation. \textit{Proposed.} \textbf{[6]}}
\end{minipage}\par\vspace{4pt}
\Panel{Ready for normal care}{Roots anchor the cutting and new growth continues. Return gradually to the established-plant care cues; propagation moisture is not a permanent watering schedule.\par Start date: \hrulefill\quad Method: \hrulefill\par First successful growth / conditions: \hrulefill}
\Sources{Evidence: 1 = KAL-RHS RHS; 2 = KAL-IA Iowa State; 3 = KAL-PROP Iowa State; 4 = KAL-HARD RHS; 6 = KAL-MSU Montana State. Proposed = proposed starting point, not a published ratio or household measurement. Full titles, URLs and scopes: entry sources.yaml and companion worksheets. Conversions rounded; source units authoritative.}{kalanchoe-desert / propagation}
